\lecture{2}{Node Embeddings}{Thursday, September 24, 2026}\label{ch:lec02}

\textbf{Outline.} Traditional ML vs.\ graph representation learning $\to$
encoder--decoder framework (shallow encoding) $\to$ random-walk embeddings
(objective, negative sampling, SGD; DeepWalk, node2vec) $\to$ embedding entire
graphs $\to$ matrix factorisation view $\to$ uses and limitations.

% =====================================================================
\section{Node Embeddings: Motivation}\label{sec:lec02-motivation}

\subsection{Traditional ML on graphs}\label{sec:lec02-traditional}
\slides{7--9}

\begin{itemize}
  \item Given an input graph, \textbf{extract node, link and graph-level
    features}, then learn a model (neural network, etc.) that maps features to
    labels:
\end{itemize}
\[
  \text{Input graph}
  \xrightarrow[\text{(node-, edge-, graph-level)}]{\text{feature engineering}}
  \text{Structured features}\to
  \underbrace{\text{Learning algorithm}\to\text{Prediction}}_{\text{downstream prediction task}}
\]

\paragraph{Example of node features: graphlets} (Pr\v{z}ulj et al.,
\emph{Bioinformatics} 2004).
\begin{itemize}
  \item \term{Graphlets}: \textbf{rooted connected induced non-isomorphic}
    subgraphs. Induced = take some nodes and \emph{all} the edges between them.
    Graphlet ID = root/``position'' of node $u$ (IDs $0$--$72$ over shapes
    $G_0,\dots,G_{29}$).
  \item \textbf{There are 73 different graphlets on up to 5 nodes.}
  \item Slide-8 example (edges $A\!-\!C$, $B\!-\!C$, $B\!-\!u$, $C\!-\!u$,
    $u\!-\!E$): $u$ has graphlets $0,1,2,3,5,10,11,\ldots$
  \item \term{Graphlet degree vector (GDV)}: a count vector of graphlets rooted
    at a given node. Example (\cref{fig:lec02-gdv}): graphlets on up to 3 nodes
    $a$ (edge), $b$ (triangle), $c$ (centre of a path), $d$ (end of a path)
    $\Rightarrow$ $\mathrm{GDV}(u)=[a,b,c,d]=[2,1,0,2]$.
\end{itemize}

\begin{figure}[htbp]
  \centering
  \begin{tikzpicture}[scale=0.85]
    \tikzset{lec02root/.style={vtx, fill=cautionred!70}}
    \node[font=\small] at (1.0,1.2) {$G$};
    \node[target] (u) at (0,0) {$u$};
    \node[vtx] (x) at (1.0,0.6) {};
    \node[vtx] (y) at (1.0,-0.6) {};
    \node[vtx] (xp) at (2.1,0.6) {};
    \node[vtx] (yp) at (2.1,-0.6) {};
    \draw[edge] (u)--(x) (u)--(y) (x)--(y) (x)--(xp) (y)--(yp);
    \node[font=\small, anchor=west] at (3.3,1.2) {root position (red): count at $u$};
    \node[lec02root] (a1) at (3.8,0.55) {}; \node[vtx] (a2) at (3.8,-0.4) {};
    \draw[edge] (a1)--(a2);
    \node[font=\small] at (3.8,-1.0) {$a$: 2};
    \node[lec02root] (b1) at (5.2,0.55) {}; \node[vtx] (b2) at (4.75,-0.4) {};
    \node[vtx] (b3) at (5.65,-0.4) {};
    \draw[edge] (b1)--(b2) (b1)--(b3) (b2)--(b3);
    \node[font=\small] at (5.2,-1.0) {$b$: 1};
    \node[lec02root] (c1) at (6.85,0.55) {}; \node[vtx] (c2) at (6.4,-0.4) {};
    \node[vtx] (c3) at (7.3,-0.4) {};
    \draw[edge] (c1)--(c2) (c1)--(c3);
    \node[font=\small] at (6.85,-1.0) {$c$: 0};
    \node[vtx] (d1) at (8.5,0.55) {}; \node[lec02root] (d2) at (8.05,-0.4) {};
    \node[vtx] (d3) at (8.95,-0.4) {};
    \draw[edge] (d1)--(d2) (d1)--(d3);
    \node[font=\small] at (8.5,-1.0) {$d$: 2};
  \end{tikzpicture}
  \caption{GDV of node $u$ from graphlets on up to 3 nodes: $[2,1,0,2]$ (slide 9).}
  \label{fig:lec02-gdv}
\end{figure}

\subsection{Graph representation learning}\label{sec:lec02-grl}
\slides{10--13}

\begin{keyidea}
  Graph Representation Learning alleviates the need to do feature engineering
  \textbf{every single time}: \emph{representation learning} automatically
  learns the features (replaces the feature-engineering step of the pipeline).
\end{keyidea}
\begin{itemize}
  \item \textbf{Goal:} efficient \textbf{task-independent} feature learning for
    machine learning with graphs. Map each node $u$ to a vector, its
    \term{feature representation / embedding}:
\end{itemize}
\begin{equation}\label{eq:lec02-f}
  f: u\to\R^d
\end{equation}
\begin{itemize}
  \item \textbf{Why embedding? Task: map nodes into an embedding space.}
    \begin{itemize}
      \item Similarity of embeddings between nodes indicates their similarity in
        the network (e.g.\ both nodes are close to each other / connected by an edge).
      \item Encode network information.
      \item Potentially used for many downstream predictions: node
        classification, link prediction, graph classification, anomalous node
        detection, clustering, \ldots
    \end{itemize}
  \item \textbf{Example:} 2D embedding of the nodes of \textbf{Zachary's Karate
    Club} network (Perozzi et al., DeepWalk, KDD 2014): input graph with
    coloured communities $\to$ output points in $\R^2$, where nodes of the same
    community lie close together.
\end{itemize}

% =====================================================================
\section{Encoder and Decoder}\label{sec:lec02-encdec}

\subsection{Setup and goal}\label{sec:lec02-setup}
\slides{15--19}

\begin{itemize}
  \item Assume an (undirected) graph $G$: $V$ is the vertex set, $\mA$ the
    adjacency matrix (assume binary). \textbf{For simplicity: no node features
    or extra information is used.} Toy example in \cref{fig:lec02-toy}.
  \item \textbf{Goal:} encode nodes so that \textbf{similarity in the embedding
    space} (e.g.\ dot product) approximates \textbf{similarity in the graph}:
\end{itemize}
\begin{equation}\label{eq:lec02-goal}
  \underbrace{\similarity(u,v)}_{\text{in the original network}}
  \;\approx\;
  \underbrace{\vz_v\tp\vz_u}_{\text{similarity of the embedding}}
\end{equation}
\begin{itemize}
  \item ``Need to define!'': both $\similarity(u,v)$ and the encoder
    $\ENC(u)=\vz_u$, $\ENC(v)=\vz_v$.
  \item \textbf{Learning node embeddings:}
    (1)~\textbf{Encoder} maps from nodes to embeddings;
    (2)~\textbf{define a node similarity function} (a measure of similarity in
    the original network);
    (3)~\textbf{Decoder} $\DEC$ maps from embeddings to the similarity score;
    (4)~\textbf{optimise the parameters of the encoder} so that
    $\similarity(u,v)\approx\DEC(\vz_v\tp\vz_u)$, here simply $\vz_v\tp\vz_u$.
  \item \textbf{Two key components:}
    \begin{itemize}
      \item \term{Encoder}: maps each node to a low-dimensional vector,
        $\ENC(v)=\vz_v$ ($v$: node in the input graph; $\vz_v$: $d$-dimensional
        embedding).
      \item \term{Similarity function}: specifies how the relationships in vector
        space map to the relationships in the original network,
        $\similarity(u,v)\approx\vz_v\tp\vz_u$; the right side, the dot product
        between node embeddings, is the \term{decoder}.
    \end{itemize}
\end{itemize}

\begin{figure}[htbp]
  \centering
  \begin{tikzpicture}[scale=0.85]
    \node[vertex] (t1) at (0,0) {1};
    \node[vertex] (t2) at (1.2,0.45) {2};
    \node[vertex] (t3) at (2.5,1.0) {3};
    \node[vertex] (t4) at (0.35,1.45) {4};
    \draw[edge] (t1)--(t2) (t1)--(t4) (t2)--(t4) (t3)--(t4);
    \node[font=\small] at (1.2,-0.7) {$V=\{1,2,3,4\}$};
    \node at (6.0,0.6) {$\mA=\begin{pmatrix}0&1&0&1\\1&0&0&1\\0&0&0&1\\1&1&1&0\end{pmatrix}$};
  \end{tikzpicture}
  \caption{Toy graph and its adjacency matrix (slide 15; reused on slides 66, 72, 74).}
  \label{fig:lec02-toy}
\end{figure}

\subsection{``Shallow'' encoding}\label{sec:lec02-shallow}
\slides{20--23}

\begin{definition}[Shallow encoder]\label{def:lec02-shallow}
  Simplest encoding approach: \textbf{the encoder is just an embedding-lookup}
  \begin{equation}\label{eq:lec02-lookup}
    \ENC(v)=\vz_v=\mZ\cdot v,
  \end{equation}
  \begin{itemize}
    \item $\mZ\in\R^{d\times|\calV|}$: \textbf{embedding matrix}, each
      \textbf{column} is a node embedding \textbf{[what we learn / optimise]}
      (slide 21: one column per node = the embedding vector for a specific node;
      \#rows = dimension/size of embeddings);
    \item $v\in\mathbb{I}^{|\calV|}$: indicator vector, all zeroes except a one in
      the column indicating node $v$.
  \end{itemize}
\end{definition}
\begin{itemize}
  \item \textbf{Each node is assigned a unique embedding vector} (i.e.\ we
    directly optimise the embedding of each node). Many methods: \textbf{DeepWalk,
    node2vec}.
  \item \textbf{Framework summary} (encoder + decoder):
    \begin{itemize}
      \item Shallow encoder: embedding lookup. Parameters to optimise: $\mZ$,
        which contains node embeddings $\vz_u$ for all nodes $u\in V$. Deep
        encoders come later with GNNs.
      \item \textbf{Decoder:} based on node similarity.
        \textbf{Objective:} maximise $\vz_v\tp\vz_u$ for node pairs $(u,v)$ that
        are \textbf{similar}.
    \end{itemize}
\end{itemize}

\subsection{How to define node similarity?}\label{sec:lec02-similarity}
\slides{24--26}

\begin{itemize}
  \item Key choice of methods is \textbf{how they define node similarity}. Should
    two nodes have a similar embedding if they \ldots\ \emph{are linked?}
    \emph{share neighbours?} \emph{have similar ``structural roles''?}
  \item Now: a node similarity definition that uses \textbf{random walks}, and
    how to optimise embeddings for such a similarity measure. (Slide 25 repeats
    the Karate Club example.)
  \item \textbf{Note on node embeddings:}
    \begin{itemize}
      \item \textbf{Unsupervised/self-supervised} way of learning node
        embeddings: \textbf{not} utilising node labels, \textbf{not} utilising
        node features. The goal is to directly estimate a set of coordinates
        (i.e.\ the embedding) of a node so that some aspect of the network
        structure (captured by $\DEC$) is preserved.
      \item The embeddings are \textbf{task independent}: not trained for a
        specific task but can be used for any task.
    \end{itemize}
\end{itemize}

% =====================================================================
\section{Random Walk Approaches for Node Embeddings}\label{sec:lec02-rw}

\subsection{Notation}\label{sec:lec02-notation}
\slides{28}

\begin{itemize}
  \item \textbf{Vector $\vz_u$}: the embedding of node $u$ (what we aim to find).
  \item \textbf{Probability $P(v\mid\vz_u)$}: our model prediction based on
    $\vz_u$, the \textbf{(predicted) probability} of visiting node $v$ on random
    walks starting from node $u$.
  \item Non-linear functions used to produce predicted probabilities:
    \textbf{softmax} turns a vector of $K$ real values (model predictions) into
    $K$ probabilities that sum to 1; \textbf{sigmoid} is an S-shaped function
    that turns real values into the range $(0,1)$:
\end{itemize}
\begin{equation}\label{eq:lec02-softmaxdef}
  S(\vz)[i]=\frac{e^{\vz[i]}}{\sum_{j=1}^{K}e^{\vz[j]}},
  \qquad
  \sigma(x)=\frac{1}{1+e^{-x}}
\end{equation}

\subsection{Random walks}\label{sec:lec02-randomwalks}
\slides{29--32}

\begin{itemize}
  \item \term{Random walk}: given a \emph{graph} and a \emph{starting point}, we
    \textbf{select a neighbour} of it at \textbf{random} and move to this
    neighbour; then select a neighbour of this point at random, and move to it,
    etc. The (random) sequence of points visited this way is a \textbf{random walk
    on the graph} (\cref{fig:lec02-randomwalk}).
\end{itemize}

\begin{figure}[htbp]
  \centering
  \begin{tikzpicture}[x=1.1cm,y=0.95cm]
    \node[vertex] (r1) at (1.0,2.75) {1};
    \node[vertex] (r2) at (2.0,3.3) {2};
    \node[vertex] (r3) at (2.65,2.3) {3};
    \node[target] (r4) at (1.65,1.5) {S};
    \node[hl] (r5) at (3.0,0.7) {5};
    \node[vertex] (r6) at (4.3,1.05) {6};
    \node[vertex] (r7) at (3.95,-0.25) {7};
    \node[hl] (r8) at (4.6,2.75) {8};
    \node[hl] (r9) at (4.45,4.3) {9};
    \node[vertex] (r10) at (5.85,4.45) {10};
    \node[hl] (r11) at (5.95,2.45) {11};
    \node[vertex] (r12) at (7.0,3.5) {12};
    \draw[edge] (r1)--(r2) (r1)--(r3) (r1)--(r4) (r2)--(r3) (r2)--(r8) (r3)--(r4)
                (r4)--(r5) (r5)--(r6) (r5)--(r7) (r6)--(r7) (r5)--(r8) (r8)--(r9)
                (r8)--(r11) (r9)--(r10) (r10)--(r11) (r10)--(r12) (r11)--(r12);
    \tikzset{lec02step/.style={arrow, cautionred, line width=1.0pt},
             lec02num/.style={font=\scriptsize\bfseries, text=cautionred, fill=white, inner sep=1pt}}
    \draw[lec02step, bend right=18] (r4) to node[lec02num, below left] {1} (r5);
    \draw[lec02step, bend right=18] (r5) to node[lec02num, right] {2} (r8);
    \draw[lec02step, bend left=35] (r8) to node[lec02num, left] {3} (r9);
    \draw[lec02step, bend left=35] (r9) to node[lec02num, right] {4} (r8);
    \draw[lec02step, bend right=22] (r8) to node[lec02num, below] {5} (r11);
  \end{tikzpicture}
  \caption{A 5-step random walk $S\to5\to8\to9\to8\to11$ from the start node $S$
  (smiley on slide 29).}
  \label{fig:lec02-randomwalk}
\end{figure}

\begin{keyidea}[Random-walk embeddings]
  $\vz_u\tp\vz_v\approx$ probability that $u$ and $v$ \textbf{co-occur on a
  random walk} over the graph.
\end{keyidea}
\begin{enumerate}
  \item \textbf{Estimate probability of visiting node $v$ on a random walk
    starting from node $u$} using some random walk strategy $R$: $P_R(v\mid u)$.
  \item \textbf{Optimise embeddings to encode these random walk statistics}:
    similarity in embedding space (here: dot product $=\cos(\theta)$) encodes
    random walk ``similarity'' (figure: angle $\theta$ between the two
    embeddings $\propto P_R(v\mid u)$).
\end{enumerate}
\textbf{Why random walks?}
\begin{enumerate}
  \item \textbf{Expressivity:} flexible stochastic definition of node similarity
    that \textbf{incorporates both local and higher-order neighbourhood
    information}. Idea: if a random walk starting from node $u$ visits $v$ with
    high probability, $u$ and $v$ are similar (high-order multi-hop information).
  \item \textbf{Efficiency:} do not need to consider all node pairs when
    training; \textbf{only need to consider pairs that co-occur on random walks}.
\end{enumerate}

\subsection{Random walk optimisation}\label{sec:lec02-objective}
\slides{33--37}

\begin{itemize}
  \item \textbf{Unsupervised feature learning.} Intuition: find embedding of nodes
    in $d$-dimensional space that preserves similarity. Idea: learn node
    embedding such that \textbf{nearby} nodes (in the network) are close together
    (in the embedding space; the slide says ``in the network''). Nearby nodes of
    $u$: $\NR(u)$, the neighbourhood of $u$ obtained by some \textbf{random walk
    strategy $R$}.
  \item \textbf{Feature learning as optimisation.} Given $G=(V,E)$, learn a
    mapping $f:u\to\R^d$, $f(u)=\vz_u$, with the \textbf{log-likelihood
    objective}
\end{itemize}
\begin{equation}\label{eq:lec02-loglik}
  \argmax_{z}\;\sum_{u\in V}\log P\big(\NR(u)\mid\vz_u\big)
  \qquad\text{(maximum likelihood objective)}
\end{equation}
\begin{itemize}
  \item Given node $u$, we want to learn feature representations that are
    predictive of the nodes in its random walk neighbourhood $\NR(u)$.
  \item \textbf{Procedure:}
    (1)~run \textbf{short fixed-length random walks} starting from each node $u$
    in the graph using some random walk strategy $R$;
    (2)~for each node $u$ collect $\NR(u)$, the \textbf{multiset} of nodes visited
    on random walks starting from $u$ (repeat elements: nodes can be visited
    multiple times);
    (3)~optimise embeddings according to: \emph{given node $u$, predict its
    neighbours $\NR(u)$}, i.e.\ \cref{eq:lec02-loglik}.
  \item \textbf{Equivalently} (with $\log P(\NR(u)\mid\vz_u)=\sum_{v\in\NR(u)}\log
    P(v\mid\vz_u)$), minimise the negative log-likelihood of the random walk
    neighbourhoods:
\end{itemize}
\begin{equation}\label{eq:lec02-nll}
  \argmin_{z}\;\Loss=\sum_{u\in V}\;\sum_{v\in\NR(u)}-\log\big(P(v\mid\vz_u)\big)
\end{equation}
\begin{itemize}
  \item \textbf{Parameterise $P(v\mid\vz_u)$ using softmax.} Why softmax? We want
    node $v$ to be most similar to node $u$ (out of all nodes $n$). Intuition:
    $\sum_i\exp(x_i)\approx\max_i\exp(x_i)$.
\end{itemize}
\begin{equation}\label{eq:lec02-softmax}
  P(v\mid\vz_u)=\frac{\exp(\vz_u\tp\vz_v)}{\sum_{n\in V}\exp(\vz_u\tp\vz_n)}
\end{equation}
\textbf{Putting it all together:}
\begin{equation}\label{eq:lec02-loss}
  \Loss=
  \underbrace{\sum_{u\in V}}_{\text{all nodes }u}\;
  \underbrace{\sum_{v\in\NR(u)}}_{\substack{\text{nodes $v$ seen on random}\\\text{walks starting from }u}}
  -\log\Bigg(
  \underbrace{\frac{\exp(\vz_u\tp\vz_v)}{\sum_{n\in V}\exp(\vz_u\tp\vz_n)}}_{\substack{\text{predicted probability of $u$ and}\\\text{$v$ co-occurring on random walk}}}
  \Bigg)
\end{equation}
\textbf{Optimising random walk embeddings $=$ finding embeddings $\vz_u$ that
minimise $\Loss$.}

\subsection{Negative sampling}\label{sec:lec02-negsampling}
\slides{38--41}

\begin{itemize}
  \item \textbf{But doing this naively is too expensive!} Nested sum over nodes
    ($\sum_{u\in V}$ and the softmax's $\sum_{n\in V}$) gives
    $\bigO(|V|^2)$ complexity. \textbf{The normalisation term from the softmax
    is the culprit} \ldots\ can we approximate it?
  \item \textbf{Solution: negative sampling.} \textbf{Instead of normalising
    w.r.t.\ all nodes, just normalise against $k$ random ``negative samples''
    $n_i$:}
\end{itemize}
\begin{equation}\label{eq:lec02-negsamp}
  \log\left(\frac{\exp(\vz_u\tp\vz_v)}{\sum_{n\in V}\exp(\vz_u\tp\vz_n)}\right)
  \approx
  \log\big(\sigma(\vz_u\tp\vz_v)\big)
  +\sum_{i=1}^{k}\log\big(\sigma(-\vz_u\tp\vz_{n_i})\big),
  \qquad n_i\sim P_V
\end{equation}
\begin{itemize}
  \item $\sigma$: sigmoid function (makes each term a ``probability'' between 0
    and 1); $P_V$: random distribution over nodes. Negative sampling allows for
    \textbf{quick likelihood calculation}.
  \item \textbf{Why is the approximation valid?} Technically, a different
    objective. But negative sampling is a form of \textbf{Noise Contrastive
    Estimation (NCE)}, which approx.\ maximises the log probability of softmax.
    The new formulation corresponds to using a \textbf{logistic regression}
    (sigmoid) to distinguish the target node $v$ from nodes $n_i$ sampled from
    the background distribution $P_V$ (arXiv:1402.3722).
  \item \textbf{Sample $k$ negative nodes $n_i$ each with probability
    proportional to its degree.}
  \item Two considerations for $k$ (\# negative samples): (1)~higher $k$ gives
    more robust estimates; (2)~higher $k$ corresponds to higher bias on negative
    events. \textbf{In practice $k=5$--$20$.}
  \item Can a negative sample be any node or only nodes not on the walk? People
    often sample \textbf{any} node (for efficiency).
\end{itemize}
\begin{caution}
  Slide 40 has $-\log(\cdot)$ on the left but $+\log\sigma$ terms on the right.
  Correct: no minus on the left (as on slide 41, \cref{eq:lec02-negsamp}).
\end{caution}

\subsection{Stochastic gradient descent}\label{sec:lec02-sgd}
\slides{42--44}

\begin{itemize}
  \item \textbf{Gradient descent} (a simple way to minimise $\Loss$): initialise
    $\vz_u$ at some randomised value for all nodes $u$; iterate until
    convergence: for all $u$ compute the derivative $\partial\Loss/\partial\vz_u$,
    then for all $u$ make a step in the reverse direction of the derivative
    ($\eta$: learning rate):
\end{itemize}
\begin{equation}\label{eq:lec02-gd}
  \vz_u\leftarrow\vz_u-\eta\,\frac{\partial\Loss}{\partial\vz_u}
\end{equation}
\begin{itemize}
  \item \textbf{Stochastic gradient descent}: instead of evaluating gradients
    over all examples, evaluate it for each \textbf{individual} training example.
    Initialise $\vz_u$ randomly for all $u$; iterate until convergence: sample a
    node $u$, for all $v$ calculate the gradient $\partial\Loss^{(u)}/\partial\vz_v$
    and update $\vz_v$:
\end{itemize}
\begin{equation}\label{eq:lec02-sgd}
  \Loss^{(u)}=\sum_{v\in\NR(u)}-\log\big(P(v\mid\vz_u)\big),
  \qquad
  \vz_v\leftarrow\vz_v-\eta\,\frac{\partial\Loss^{(u)}}{\partial\vz_v}
\end{equation}
\paragraph{Random walks: summary.}
(1)~Run \textbf{short fixed-length} random walks starting from each node on the
graph. (2)~For each node $u$ collect $\NR(u)$, the multiset of nodes visited on
random walks starting from $u$. (3)~Optimise embeddings $\mZ$ using SGD on
$\Loss=\sum_{u\in V}\sum_{v\in\NR(u)}-\log(P(v\mid\vz_u))$ ---
\emph{we can efficiently approximate this using negative sampling!}

\subsection{How should we randomly walk? DeepWalk and node2vec}\label{sec:lec02-node2vec}
\slides{45--49}

\begin{itemize}
  \item So far: how to optimise embeddings \emph{given} a random walk strategy
    $R$. \textbf{What strategies should we use to run these random walks?}
  \item Simplest idea: \textbf{just run fixed-length, unbiased random walks
    starting from each node} (\term{DeepWalk}, Perozzi et al., KDD). Issue: such
    notion of similarity is \textbf{too constrained}. How can we generalise this?
  \item \textbf{Overview of \term{node2vec}} (Grover et al., KDD):
    \begin{itemize}
      \item \textbf{Goal:} embed nodes with similar network neighbourhoods close
        in the feature space; framed as a maximum likelihood optimisation
        problem, independent of the downstream prediction task.
      \item \textbf{Key observation:} a flexible notion of network neighbourhood
        $\NR(u)$ of node $u$ leads to rich node embeddings.
      \item Develop a \textbf{biased 2nd order random walk} $R$ to generate the
        network neighbourhood $\NR(u)$ of node $u$.
    \end{itemize}
  \item \textbf{Biased walks.} Idea: use flexible, biased random walks that can
    trade off between \textbf{local} and \textbf{global} views of the network.
    Two classic strategies to define a neighbourhood $\NR(u)$
    (\cref{fig:lec02-bfsdfs}); walk of length 3 ($\NR(u)$ of size 3):
\end{itemize}
\begin{align}
  N_{BFS}(u)&=\{s_1,s_2,s_3\} &&\text{\textbf{Local} microscopic view}\label{eq:lec02-bfs}\\
  N_{DFS}(u)&=\{s_4,s_5,s_6\} &&\text{\textbf{Global} macroscopic view}\label{eq:lec02-dfs}
\end{align}
\begin{itemize}
  \item \textbf{BFS vs.\ DFS:} BFS: $\NR(\cdot)$ provides a \textbf{micro-view}
    of the neighbourhood; DFS: $\NR(\cdot)$ provides a \textbf{macro-view} of the
    neighbourhood.
\end{itemize}

\begin{figure}[htbp]
  \centering
  \begin{tikzpicture}[scale=0.95]
    \tikzset{lec02bfsn/.style={vertex, fill=cautionred!20},
             lec02dfsn/.style={vertex, fill=keyblue!20}}
    \node[lec02bfsn] (s1) at (0.73,2.07) {$s_1$};
    \node[lec02bfsn] (s2) at (2.74,2.0) {$s_2$};
    \node[vertex, fill=black!25] (u) at (1.81,1.07) {$u$};
    \node[lec02bfsn] (s3) at (0.67,0.1) {$s_3$};
    \node[lec02dfsn] (s4) at (2.72,0.26) {$s_4$};
    \node[lec02dfsn] (s5) at (4.2,0.15) {$s_5$};
    \node[vertex] (s7) at (4.13,1.85) {$s_7$};
    \node[lec02dfsn] (s6) at (4.86,1.04) {$s_6$};
    \node[vertex] (s8) at (5.6,1.97) {$s_8$};
    \node[vertex] (s9) at (5.74,0.36) {$s_9$};
    \draw[edge] (s1)--(s2) (s1)--(u) (s1)--(s3) (u)--(s2) (u)--(s3) (u)--(s4)
                (s2)--(s4) (s2)--(s5) (s3)--(s4) (s4)--(s5) (s5)--(s6) (s5)--(s7)
                (s6)--(s7) (s7)--(s8) (s6)--(s8) (s6)--(s9) (s8)--(s9);
    \tikzset{lec02bfsa/.style={arrow, cautionred, line width=1.1pt, shorten <=2pt, shorten >=2pt},
             lec02dfsa/.style={arrow, keyblue, line width=1.1pt, shorten <=2pt, shorten >=2pt}}
    \draw[lec02bfsa, bend left=12] (u) to (s1);
    \draw[lec02bfsa, bend left=12] (u) to (s2);
    \draw[lec02bfsa, bend left=12] (u) to (s3);
    \draw[lec02dfsa, bend left=14] (u) to (s4);
    \draw[lec02dfsa, bend left=14] (s4) to (s5);
    \draw[lec02dfsa, bend left=14] (s5) to (s6);
    \draw[arrow, cautionred, line width=1.1pt] (6.5,1.5) -- (7.1,1.5) node[right, text=black, font=\small] {BFS};
    \draw[arrow, keyblue, line width=1.1pt] (6.5,0.8) -- (7.1,0.8) node[right, text=black, font=\small] {DFS};
  \end{tikzpicture}
  \caption{BFS vs.\ DFS neighbourhoods of $u$ of size 3 (slides 47--48).}
  \label{fig:lec02-bfsdfs}
\end{figure}

\subsection{The biased random walk: parameters \texorpdfstring{$p$}{p} and \texorpdfstring{$q$}{q}}\label{sec:lec02-biased}
\slides{50--53}

\begin{figure}[htbp]
  \centering
  \begin{minipage}{0.55\textwidth}
    \centering
    \begin{tikzpicture}[scale=1.05]
      \node[vertex, fill=black!15] (u) at (0,0) {$u$};
      \node[vertex, fill=YellowOrange!40] (t) at (1.7,0.08) {$t$};
      \node[vertex, fill=cautionred!25] (s1) at (1.75,1.45) {$s_1$};
      \node[vertex, fill=connectgreen!40] (w) at (3.0,0.85) {$w$};
      \node[vertex, fill=YellowOrange!20] (s2) at (4.5,1.4) {$s_2$};
      \node[vertex, fill=YellowOrange!20] (s3) at (4.4,0.0) {$s_3$};
      \draw[edge] (u)--(t) (t)--(s1) (t)--(w) (s1)--(w) (w)--(s2) (w)--(s3);
      \draw[edge] (t)--(1.45,-0.55);
      \tikzset{lec02tr/.style={arrow, line width=1.0pt, shorten <=2pt, shorten >=2pt}}
      \draw[lec02tr, keyblue, bend left=30] (w) to (t);
      \node[font=\small, text=keyblue] at (2.75,0.1) {$1/p$};
      \draw[lec02tr, cautionred, bend right=30] (w) to node[above right=-1pt, font=\small] {$1$} (s1);
      \draw[lec02tr, Brown, bend left=18] (w) to node[above, font=\small] {$1/q$} (s2);
      \draw[lec02tr, Brown, bend left=18] (w) to node[above right=-1pt, font=\small] {$1/q$} (s3);
    \end{tikzpicture}
  \end{minipage}%
  \begin{minipage}{0.4\textwidth}
    \centering\small\setlength{\tabcolsep}{4pt}
    \begin{tabular}{@{}ccc@{}}
      \toprule
      Target $x$ & Prob.\ (unnorm.) & Dist.\ $(t,x)$\\
      \midrule
      $t$   & $1/p$ & 0\\
      $s_1$ & $1$   & 1\\
      $s_2$ & $1/q$ & 2\\
      $s_3$ & $1/q$ & 2\\
      \bottomrule
    \end{tabular}
  \end{minipage}
  \caption{One step of the biased walk (slides 51--53): walker came over $(t,w)$
  and is at $w$; unnormalised transition probabilities segmented by distance
  from $t$.}
  \label{fig:lec02-n2vstep}
\end{figure}

\begin{itemize}
  \item \textbf{Interpolating BFS and DFS:} biased fixed-length random walk $R$
    that, given a node $u$, generates neighbourhood $\NR(u)$. Two parameters:
    \begin{itemize}
      \item \term{Return parameter $p$}: return back to the previous node.
      \item \term{In-out parameter $q$}: moving outwards (DFS) vs.\ inwards (BFS)
        from the previous node; intuitively, $q$ is the ``ratio'' of BFS vs.\ DFS.
    \end{itemize}
    Specify how a \textbf{single step} is performed; the random walk is then just
    a sequence of these steps.
  \item \textbf{One step:} define the walk by the transition probabilities on
    edges adjacent to the current node $w$. The walk \textbf{just traversed edge
    $(t,w)$} and \textbf{is now at $w$}; we specify edge transition probs.\ out of
    $w$. \textbf{Insight:} neighbours of $w$ can only be at distance 0, 1 or 2
    from $t$.
  \item $p,q$ model transition probabilities ($p$: return parameter; $q$:
    ``walk away'' parameter); unnormalised probabilities for $w\to x$
    (\cref{fig:lec02-n2vstep}):
\end{itemize}
\begin{equation}\label{eq:lec02-alpha}
  \alpha_{pq}(t,x)=
  \begin{cases}
    1/p & \text{if } d(t,x)=0 \quad(x=t),\\
    1   & \text{if } d(t,x)=1,\\
    1/q & \text{if } d(t,x)=2.
  \end{cases}
\end{equation}
\begin{itemize}
  \item Normalise over the neighbours of $w$ (in the example: divide by
    $1/p+1+2/q$).
  \item \textbf{BFS-like walk: low value of $p$.} \textbf{DFS-like walk: low
    value of $q$.} $\NR(u)$ are the nodes visited by the biased walk.
\end{itemize}

\begin{caution}
  Slide 53 says ``walker came over edge $(s_1,w)$''; it should be $(t,w)$ (the
  table's distances are measured from $t$).
\end{caution}

\subsection{node2vec algorithm and other random walk ideas}\label{sec:lec02-n2v-alg}
\slides{54--57}

\begin{itemize}
  \item \textbf{node2vec algorithm:}
    (1)~compute edge transition probabilities: for each edge $(t,w)$ compute edge
    walk probabilities (based on $p,q$) of edges $(w,\cdot)$;
    (2)~simulate $r$ random walks of length $l$ starting from each node $u$;
    (3)~optimise the node2vec objective using SGD.
    \textbf{Linear-time} complexity; all 3 steps are \textbf{individually
    parallelisable}.
  \item \textbf{Other random walk ideas:}
    \begin{itemize}
      \item Different kinds of biased random walks: based on node attributes
        (Dong et al.); based on learned weights (Abu-El-Haija et al.).
      \item Alternative optimisation schemes: directly optimise based on 1-hop and
        2-hop random walk probabilities (as in LINE from Tang et al.).
      \item Network preprocessing techniques: run random walks on modified
        versions of the original network (e.g.\ Ribeiro et al.\ struct2vec,
        Chen et al.\ HARP).
    \end{itemize}
  \item \textbf{Summary so far.} \textbf{Core idea:} embed nodes so that distances
    in embedding space reflect node similarities in the original network.
    Different notions of node similarity: na\"ive (similar if two nodes are
    connected); random walk approaches (covered today).
  \item \textbf{So, what method should I use?} No one method wins in all cases;
    e.g.\ node2vec performs better on node classification while alternative
    methods perform better on link prediction (Goyal and Ferrara, survey).
    Random walk approaches are generally more efficient. \textbf{In general:} must
    choose the definition of node similarity that matches your application.
\end{itemize}

% =====================================================================
\section{Embedding Entire Graphs}\label{sec:lec02-graphemb}
\slides{59--63}

\begin{itemize}
  \item \textbf{Goal:} embed a subgraph or an entire graph $G$; graph embedding
    $\vz_G$. \textbf{Tasks:} classifying toxic vs.\ non-toxic molecules;
    identifying anomalous graphs.
  \item \textbf{Approach 1} (simple but effective): run a standard graph
    embedding technique \emph{on} the (sub)graph $G$, then just \textbf{sum (or
    average)} the node embeddings in the (sub)graph $G$. Used by Duvenaud et al.\
    to classify molecules based on their graph structure.
\end{itemize}
\begin{equation}\label{eq:lec02-sumpool}
  \vz_G=\sum_{v\in G}\vz_v
\end{equation}
\begin{itemize}
  \item \textbf{Approach 2:} introduce a \textbf{``virtual node''} to represent
    the (sub)graph $S$ (figure: an extra node connected to all nodes of the
    circled subgraph $S$) and run a standard graph embedding technique; the
    virtual node's embedding is $\vz_S$. Proposed by Li et al.\ as a
    general technique for subgraph embedding.
  \item \textbf{Summary:} Approach 1: embed nodes and sum/avg them. Approach 2:
    create a super-node that spans the (sub)graph and then embed that node.
  \item \textbf{Preview -- hierarchical embeddings (\term{DiffPool}):} we can also
    \textbf{hierarchically} cluster nodes in graphs, and \textbf{sum/avg} the node
    embeddings according to these clusters. Figure: each cluster becomes one
    node of the next level: original network (5 clusters) $\to$ pooled network at
    level 1 (5 nodes, 2 clusters) $\to$ level 2 (2 nodes) $\to$ level 3 (1 node)
    $\to$ its embedding is fed to a neural network for graph classification.
\end{itemize}

% =====================================================================
\section{Matrix Factorization and Node Embeddings}\label{sec:lec02-mf}

\subsection{Edge-based similarity}\label{sec:lec02-mf-adj}
\slides{65--67}

\begin{itemize}
  \item \textbf{Recall:} encoder as an embedding lookup ($\mZ$: one column per
    node). Objective: maximise $\vz_v\tp\vz_u$ for node pairs $(u,v)$ that are
    similar.
  \item Simplest \textbf{node similarity}: nodes $u,v$ are similar if they are
    connected by an edge. This means $\vz_v\tp\vz_u=A_{u,v}$, the $(u,v)$ entry of
    the adjacency matrix $\mA$. Figure (toy graph of \cref{fig:lec02-toy}):
    highlighted entry $A_{2,4}=1$ (edge $2$--$4$) $=$ row $u$ of $\mZ\tp$
    ($\vz_u$) $\times$ column $v$ of $\mZ$ ($\vz_v$). Therefore:
\end{itemize}
\begin{equation}\label{eq:lec02-ztz}
  \vz_v\tp\vz_u=A_{u,v}
  \quad\Longrightarrow\quad
  \mZ\tp\mZ=\mA
\end{equation}
\begin{itemize}
  \item The embedding dimension $d$ (number of rows in $\mZ$) is much smaller than
    the number of nodes $n$ $\Rightarrow$ exact factorisation $\mA=\mZ\tp\mZ$ is
    generally not possible. However, we can learn $\mZ$ approximately:
\end{itemize}
\begin{equation}\label{eq:lec02-mf-obj}
  \min_{\mZ}\;\big\lVert\mA-\mZ\tp\mZ\big\rVert_2
\end{equation}
\begin{itemize}
  \item We optimise $\mZ$ such that it minimises the L2 norm (Frobenius norm) of
    $\mA-\mZ\tp\mZ$. Note: today we used softmax instead of L2, but the goal to
    approximate $\mA$ with $\mZ\tp\mZ$ is the same.
\end{itemize}
\begin{keyidea}
  Conclusion: \textbf{inner product decoder with node similarity defined by edge
  connectivity is equivalent to matrix factorisation of $\mA$.}
\end{keyidea}

\subsection{Random walk-based similarity}\label{sec:lec02-mf-rw}
\slides{68--69}

\begin{itemize}
  \item DeepWalk and node2vec have a more complex \textbf{node similarity}
    definition based on random walks.
  \item \textbf{DeepWalk is equivalent to matrix factorisation} of the following
    complex matrix expression (Qiu et al., \emph{Network Embedding as Matrix
    Factorization: Unifying DeepWalk, LINE, PTE, and node2vec}):
\end{itemize}
\begin{equation}\label{eq:lec02-netmf}
  \log\!\left(\mathrm{vol}(G)\left(\frac{1}{T}\sum_{r=1}^{T}\big(\mD^{-1}\mA\big)^{r}\right)\mD^{-1}\right)-\log b,
  \qquad
  \mathrm{vol}(G)=\sum_i\sum_jA_{i,j}
\end{equation}
\begin{itemize}
  \item $\mathrm{vol}(G)$: \textbf{volume of graph}; $\mD$: \textbf{diagonal
    matrix}, $D_{u,u}=\deg(u)$; $(\mD^{-1}\mA)^r$: \textbf{power of normalised
    adjacency matrix}; $T$: \textbf{context window size}, $T=|\NR(u)|$ (see Lec 3
    slide 30); $b$:
    \textbf{number of negative samples}.
  \item \textbf{node2vec} can also be formulated as a matrix factorisation
    (albeit a more complex matrix); refer to the paper for details.
\end{itemize}

% =====================================================================
\section{Using Embeddings, Summary and Limitations}\label{sec:lec02-use}

\subsection{How to use embeddings}\label{sec:lec02-howtouse}
\slides{70--71}

How to use embeddings $\vz_i$ of nodes:
\begin{itemize}
  \item \textbf{Clustering/community detection:} cluster points $\vz_i$.
  \item \textbf{Node classification:} predict label of node $i$ based on $\vz_i$.
  \item \textbf{Link prediction:} predict edge $(i,j)$ based on $(\vz_i,\vz_j)$;
    we can concatenate, avg, product, or take a difference between the
    embeddings ($*$ = Hadamard, per coordinate product):
\end{itemize}
\begin{align}
  \text{Concatenate: }& f(\vz_i,\vz_j)=g([\vz_i,\vz_j]) &
  \text{Hadamard: }& f(\vz_i,\vz_j)=g(\vz_i*\vz_j)\notag\\
  \text{Sum/Avg: }& f(\vz_i,\vz_j)=g(\vz_i+\vz_j) &
  \text{Distance: }& f(\vz_i,\vz_j)=g(\norm{\vz_i-\vz_j}_2)\label{eq:lec02-linkpred}
\end{align}
\begin{itemize}
  \item \textbf{Graph classification:} graph embedding $\vz_G$ via aggregating
    node embeddings or virtual-node; predict label based on $\vz_G$.
  \item \textbf{Today's summary:} graph representation learning = a way to learn
    \textbf{node and graph embeddings} for downstream tasks \textbf{without
    feature engineering}. Encoder-decoder framework (encoder: embedding lookup;
    decoder: predict score based on embedding to match node similarity). Node
    similarity measure: (biased) random walk (examples: DeepWalk, node2vec).
    Extension to graph embedding: node embedding aggregation.
\end{itemize}

\subsection{Limitations}\label{sec:lec02-limitations}
\slides{72--74}

Limitations of node embeddings via matrix factorisation and random walks:
\begin{enumerate}
  \item \textbf{Transductive (not inductive) method:} cannot obtain embeddings for
    nodes not in the training set; cannot apply to new graphs. If you apply
    DeepWalk to the same graph multiple times, you'll get a different embedding
    each time. Why? (stochastic: random walks, initialisation, SGD).
    Example: a newly added node 5 at test time (e.g.\ new user in a social
    network), attached to node 2 of the toy training graph: cannot compute its
    embedding with DeepWalk/node2vec $\Rightarrow$ need to recompute all node
    embeddings.
  \item \textbf{Cannot capture structural similarity} (\cref{fig:lec02-struct}):
    nodes 1 and 11 are \textbf{structurally similar} (part of one triangle,
    degree 2, \ldots) but have very \textbf{different} embeddings: it's unlikely
    that a random walk will reach node 11 from node 1. \textbf{DeepWalk and
    node2vec do not capture structural similarity.}
  \item \textbf{Cannot utilise node, edge and graph features}, e.g.\ a feature
    vector per node (protein properties in a protein-protein interaction graph):
    DeepWalk/node2vec embeddings do not incorporate such node features.
\end{enumerate}

\begin{figure}[htbp]
  \centering
  \begin{tikzpicture}[scale=0.8]
    \node[target] (g1) at (0,0) {1};
    \node[vertex] (g2) at (1.3,0.35) {2};
    \node[vertex] (g3) at (0.4,1.2) {3};
    \node[vertex] (g4) at (2.3,0.95) {4};
    \node[vertex] (g5) at (3.4,0.95) {5};
    \node[vertex] (g10) at (7.0,0.95) {10};
    \node[vertex] (g13) at (8.2,1.35) {13};
    \node[vertex] (g12) at (9.3,0.5) {12};
    \node[target] (g11) at (7.95,0.0) {11};
    \draw[edge] (g1)--(g2) (g1)--(g3) (g2)--(g3) (g3)--(g4) (g4)--(g5)
                (g10)--(g13) (g13)--(g12) (g12)--(g11) (g11)--(g13);
    \draw[edge, dashed] (g5)--(g10);
  \end{tikzpicture}
  \caption{Nodes 1 and 11 are structurally similar but far apart (slide 73).}
  \label{fig:lec02-struct}
\end{figure}

\begin{keyidea}
  Solution to these limitations: \textbf{Deep Representation Learning and Graph
  Neural Networks} (covered in depth next).
\end{keyidea}

% ---------------------------------------------------------------------
\bigskip\phantomsection
\addcontentsline{toc}{section}{Lecture 2 at a glance}
\begin{summarybox}[Lecture 2 at a glance]
  \begin{itemize}
    \item \textbf{Representation learning} replaces per-task feature engineering:
      learn $f:u\to\R^d$, task-independent, unsupervised (no labels, no features).
    \item \textbf{Encoder--decoder:} $\ENC(v)=\vz_v$;
      $\similarity(u,v)\approx\vz_v\tp\vz_u$ (decoder = dot product). Key choice:
      how to define node similarity.
    \item \textbf{Shallow encoder:} $\ENC(v)=\vz_v=\mZ\cdot v$,
      $\mZ\in\R^{d\times|\calV|}$ (one column per node; DeepWalk, node2vec).
    \item \textbf{Random-walk similarity:} $\vz_u\tp\vz_v\approx$ prob.\ that $u,v$
      co-occur on a walk; expressive (multi-hop) and efficient (co-occurring pairs only).
    \item \textbf{Objective:} minimise
      $\Loss=\sum_{u\in V}\sum_{v\in\NR(u)}-\log P(v\mid\vz_u)$
      ($\Leftrightarrow$ maximise the log-likelihood
      $\sum_{u\in V}\log P(\NR(u)\mid\vz_u)$).
    \item \textbf{Softmax:}
      $P(v\mid\vz_u)=\exp(\vz_u\tp\vz_v)/\sum_{n\in V}\exp(\vz_u\tp\vz_n)$
      $\Rightarrow$ naive cost $\bigO(|V|^2)$.
    \item \textbf{Negative sampling:} $\log P(v\mid\vz_u)\approx\log\sigma(\vz_u\tp\vz_v)
      +\sum_{i=1}^k\log\sigma(-\vz_u\tp\vz_{n_i})$, $n_i\sim P_V\propto$ degree;
      $k=5$--$20$. Optimise with SGD:
      $\vz_v\leftarrow\vz_v-\eta\,\partial\Loss^{(u)}/\partial\vz_v$.
    \item \textbf{DeepWalk:} unbiased fixed-length walks. \textbf{node2vec:} biased
      2nd-order walk, weights $1/p,1,1/q$ for $d(t,x)=0,1,2$; BFS (local, micro)
      vs.\ DFS (global, macro); linear time, parallelisable.
    \item \textbf{Graph embedding:} $\vz_G=\sum_{v\in G}\vz_v$, virtual node,
      DiffPool (hierarchical).
    \item \textbf{Matrix factorisation:} edge similarity $\Leftrightarrow$
      $\min_{\mZ}\lVert\mA-\mZ\tp\mZ\rVert_2$; DeepWalk factorises
      $\log\big(\mathrm{vol}(G)(\frac1T\sum_{r=1}^T(\mD^{-1}\mA)^r)\mD^{-1}\big)-\log b$.
    \item \textbf{Uses:} clustering, node classification, link prediction
      (concat / Hadamard / sum / distance), graph classification.
      \textbf{Limitations:} transductive, no structural similarity, no features
      $\Rightarrow$ GNNs.
  \end{itemize}
\end{summarybox}
